% !TEX root = main.tex

% =============================
\documentclass[12pt,a4paper]{report}

% --- Font (Times New Roman asli) ---
\usepackage{fontspec}
\setmainfont{Times New Roman}[Ligatures=TeX,Renderer=Basic]

\usepackage[american,indonesian]{babel}

\usepackage{csquotes}
\usepackage[style=apa,backend=biber]{biblatex}
\addbibresource{references.bib}

% --- MAP bahasa indonesian ke mapping 'american-apa' yang tersedia ---
\DeclareLanguageMapping{indonesian}{american-apa}

% --- Override label bibliografi ke Bahasa Indonesia ---
\DefineBibliographyStrings{indonesian}{%
  bibliography = {DAFTAR PUSTAKA},
  references = {DAFTAR PUSTAKA},
  and          = {\&},
}

% Pilih bahasa tampilan dokumen jadi indonesian
\AtBeginDocument{\selectlanguage{indonesian}}

% --- Margin ---
\usepackage{geometry}
\geometry{
  a4paper,
  left=4cm,
  right=3cm,
  top=4cm,
  bottom=3cm
}

% --- Spasi 1.5 ---
\usepackage{setspace}
\onehalfspacing

% --- Paragraf (indentasi) ---
\usepackage{indentfirst}
\setlength{\parindent}{1cm}
\setlength{\parskip}{0pt}

% --- Matematika & fisika ---
\usepackage{amsmath,amssymb,amsfonts,mathtools}
\usepackage{physics}
\usepackage{bm}
\usepackage{tensor}
\usepackage{siunitx}

% --- Gambar & grafik ---
\usepackage{graphicx}
\usepackage{caption}
\usepackage{subcaption}
\usepackage{float}

% --- Tabel ---
\usepackage{booktabs}
\usepackage{multirow}
\usepackage{array}

% --- Listing kode Python ---
\usepackage{listings}
\usepackage{xcolor}

\definecolor{codebg}{RGB}{245,245,245}
\definecolor{codegreen}{RGB}{0,128,0}
\definecolor{codegray}{RGB}{128,128,128}
\definecolor{codepurple}{RGB}{128,0,128}

\lstdefinestyle{pythonstyle}{
  backgroundcolor=\color{codebg},
  basicstyle=\ttfamily\small,
  breaklines=true,
  captionpos=b,
  commentstyle=\color{codegreen},
  keywordstyle=\color{codepurple}\bfseries,
  stringstyle=\color{red!70!black},
  numberstyle=\tiny\color{codegray},
  numbers=left,
  numbersep=5pt,
  frame=single,
  rulecolor=\color{black!30},
  tabsize=4,
  language=Python,
  showstringspaces=false,
}
\lstset{style=pythonstyle}

% --- Hyperlink & referensi ---
\usepackage[hidelinks]{hyperref}
\usepackage{cleveref}

% --- Penomoran persamaan per bab ---
\numberwithin{equation}{chapter}

% --- Enumerate custom ---
\usepackage{enumitem}

% --- Format judul bab ---
\usepackage{titlesec}

\renewcommand{\thechapter}{\arabic{chapter}}
\renewcommand{\thesection}{\thechapter.\arabic{section}}
\renewcommand{\thesubsection}{\thesection.\arabic{subsection}}

\titleformat{\chapter}[display]
  {\normalfont\centering\bfseries\normalsize}
  {\MakeUppercase{BAB \Roman{chapter}}}
  {0pt}
  {\vspace{1ex}\MakeUppercase}
\titlespacing*{\chapter}{0pt}{0pt}{20pt}

\titleformat{\section}
  {\normalfont\bfseries\normalsize}
  {\thesection}{1em}{}
\titlespacing*{\section}{0pt}{12pt}{6pt}

\titleformat{\subsection}
  {\normalfont\bfseries\normalsize}
  {\thesubsection}{1em}{}
\titlespacing*{\subsection}{0pt}{10pt}{4pt}

\usepackage[none]{hyphenat}
\sloppy
